%%%%%%%%%%%%%%%%%%%%%%%%%%%%%%%%%%%%%%%%%
% Standalone ontology proposal
% Simplified PV array closed-loop control
%%%%%%%%%%%%%%%%%%%%%%%%%%%%%%%%%%%%%%%%%

\documentclass[
	letterpaper,
	10pt,
	unnumberedsections,
	twoside,
]{LTJournalArticle}

\addbibresource{patchAliasing.bib}

\usepackage{float}
\usepackage{tikz}
\usetikzlibrary{arrows.meta,positioning,shapes.geometric}
\setlength{\thanksmarkwidth}{5pt}

\runninghead{PV Array Ontology}

\setcounter{page}{1}

\title{Proposed Simplified Ontology\\for PV Array Closed-Loop Control}

\author{
	Matteo Boniotti, Gianluca Brignoli and Federico Sabbadini\thanks{This work is licensed under the Creative Commons Attribution 4.0 International License. Concurrently, all associated source code, probing scripts, and software artifacts generated throughout the project lifecycle are released under the open-source MIT license. \newline Copyright for any third-party material included in this work remains with its respective owners and must be respected accordingly. \newline \newline During the preparation of this work, the authors used AI-based tools, including GPT-5.5, Gemini 3.6 Flash, Sonnet 5, and Opus 4.6, to support language refinement and improve clarity. All generated or suggested content was carefully reviewed, revised where necessary, and validated by the authors, who assume full responsibility for the final content of this work.}
}

\hypersetup{
	pdftitle={Proposed Simplified Ontology for PV Array Closed-Loop Control},
	pdfauthor={Matteo Boniotti, Gianluca Brignoli and Federico Sabbadini},
}

\newcommand{\pv}[1]{\texttt{pv:#1}}
\newcommand{\valueof}[1]{\texttt{pv:#1}}

\begin{document}

\maketitle

This document defines a compact ontology for a photovoltaic system governed by one control agent. It connects environmental and electrical observations to anomaly recognition, Chronos-style patched representations, structural-aliasing evidence, and bounded power-routing decisions~\cite{ansariChronosLearningLanguage2024}. Routing is included only as the operational endpoint of the semantic model; the principal concern remains whether the representation preserves the information on which the agent relies.

\section{Domain Description}

The modeled system contains four energy assets: a photovoltaic array, battery storage, an aggregate internal load, and an external grid. The internal load is a passive sink representing a factory or a fleet of machines. Its demand is observed, but it has no agent, machine-level taxonomy, curtailment command, or scheduling model. The external grid is an available bidirectional boundary bus. It supplies residual demand and accepts residual production.

Exactly one \pv{ControlAgent} coordinates the four passive assets. It observes one synchronized control cycle, assigns environmental, anomaly, and observability states, selects a control mode, and emits one or more directional decisions. There is no peer-to-peer negotiation among assets.

The seven admissible transfer directions are
\begin{equation}
\label{eq:routes}
\begin{aligned}
\mathrm{PV}&\rightarrow\mathrm{Load}, &
\mathrm{PV}&\rightarrow\mathrm{Battery}, &
\mathrm{PV}&\rightarrow\mathrm{Grid},\\
\mathrm{Battery}&\rightarrow\mathrm{Load}, &
\mathrm{Battery}&\rightarrow\mathrm{Grid}, &
\mathrm{Grid}&\rightarrow\mathrm{Load},\\
&&\mathrm{Grid}&\rightarrow\mathrm{Battery}.&&
\end{aligned}
\end{equation}
They satisfy the cycle-level balance
\begin{equation}
\label{eq:power-balance}
P_{\mathrm{pv}}+P_{\mathrm{bat,dis}}+P_{\mathrm{grid,in}}
=P_{\mathrm{load}}+P_{\mathrm{bat,ch}}+P_{\mathrm{grid,out}}.
\end{equation}
All terms denote non-negative magnitudes; import and export are therefore kept separate rather than encoded by an ambiguous sign.

\section{Purpose and Scope}

The ontology provides the smallest semantic layer needed to represent the physical boundary, validated observations, environmental context, three selected anomalies, and the effect of patching on model observability. It does not model individual machines, market prices, dispatch optimisation, grid faults, or a deployed real-time controller. Load demand, battery SoC, asset limits, and the optional grid-exchange reference are asserted or simulated because they are not present in the photovoltaic telemetry dataset.

The environmental states are operational generation proxies, not meteorological diagnoses. Likewise, the routing rules describe a deterministic conceptual policy rather than the central experimental contribution. Structural aliasing remains an empirically tested property of a signal-representation pair; no fixed synthetic frequency is transferred to the one-minute photovoltaic signal.

\section{Ontology Representation}

\subsection{Core classes and closed vocabularies}

\autoref{tab:classes} defines the compact class interface. The four asset subclasses form a disjoint union; each \pv{PVSystem} contains exactly one asset of each type and is managed by exactly one control agent. Closed operational vocabularies are represented as named individuals rather than deeper class hierarchies.

\begin{table}[H]
\centering
\small
\begin{tabular}{@{}L{0.29\linewidth}L{0.64\linewidth}@{}}
\toprule
\textbf{Class} & \textbf{Role} \\
\midrule
\pv{PVSystem} & Aggregate system and scope of each control cycle. \\
\pv{EnergyAsset} & Parent of the four passive physical assets. \\
\pv{PhotovoltaicArray} & Intermittent renewable source measured through PV telemetry. \\
\pv{BatteryStorage} & Bounded storage asset that can charge, discharge, or hold. \\
\pv{InternalLoad} & Aggregate passive demand of a factory or machine fleet. \\
\pv{ExternalGrid} & Reliable bidirectional boundary bus for residual balancing. \\
\pv{ControlAgent} & The single decision-making entity. \\
\pv{SystemObservation} & Time-stamped snapshot of PV, environment, battery, load, and grid reference. \\
\pv{TimeSeriesSignal} & Physical signal and its temporal descriptors. \\
\pv{PatchedRepresentation} & A model input representation defined by patch and stride. \\
\pv{ObservabilityAssessment} & Evidence linking a signal-representation pair to an observability state. \\
\pv{GenerationContext} & Operational proxy for expected solar generation. \\
\pv{AnomalyState} & Detected physical or sensing abnormality. \\
\pv{ObservabilityState} & Semantic assessment of information retained after patching. \\
\pv{ControlMode} & Normal, conservative, or safe operating policy. \\
\pv{PowerRoute} & Closed vocabulary of the seven admissible source--destination pairs. \\
\pv{PowerRoutingDecision} & One directional transfer or hold decision with an optional bounded setpoint. \\
\bottomrule
\end{tabular}
\caption{Core classes in the simplified ontology.}
\label{tab:classes}
\end{table}

The closed vocabularies are:
\begin{description}
	\item[Generation context:] \valueof{ClearSky}, \valueof{PartialCloud}, \valueof{Overcast}, and \valueof{Nighttime}.
	\item[Anomaly:] \valueof{InverterTrip}, \valueof{RampEvent}, and \valueof{SensorMalfunction}.
	\item[Observability:] \valueof{Observable}, \valueof{PartiallyObservable}, and \valueof{StructurallyAliased}.
	\item[Control mode:] \valueof{NormalMode}, \valueof{ConservativeMode}, and \valueof{SafeMode}.
	\item[Power route:] \valueof{PVToInternalLoad}, \valueof{PVToBattery}, and \valueof{PVToExternalGrid};\newline
	\valueof{BatteryToInternalLoad} and \valueof{BatteryToExternalGrid};\newline
	\valueof{ExternalGridToInternalLoad} and \valueof{ExternalGridToBattery}.
\end{description}
Formally, each vocabulary class is an OWL-style enumeration of the listed named individuals (\texttt{oneOf}). The five vocabulary classes are pairwise disjoint, and all listed values are declared globally different (\texttt{AllDifferent}). This prevents a value such as \valueof{SafeMode} from being silently identified with \valueof{ClearSky} or \valueof{NormalMode} under open-world semantics.
Each valid time-stamped observation has exactly one generation context. Each completed observability assessment has exactly one observability state, and each routing decision records exactly one control mode. Co-issued decisions based on the same synchronized observation must record the same mode. Anomalies are not declared mutually exclusive because, for example, an abrupt inverter trip also creates a large power change. Deterministic control is obtained through the priority rules in the closed-loop section.

\subsection{Object properties}

The semantic relations in \autoref{tab:objects} connect the physical boundary to the inference and control layers. Each active \pv{PowerRoutingDecision} has exactly one route value, one source, and one distinct destination. Its non-negative setpoint magnitude is optional because \texttt{Balance} may denote passive residual exchange. A cycle that simultaneously serves the load, charges the battery, and exports a residual uses three instances of the same decision class, all based on the same observation and mode.

\begin{table}[H]
\centering
\small
\begin{tabular}{@{}L{0.25\linewidth}L{0.29\linewidth}L{0.37\linewidth}@{}}
\toprule
\textbf{Property} & \textbf{Domain} & \textbf{Range or meaning} \\
\midrule
\pv{hasAsset} & \pv{PVSystem} & \pv{EnergyAsset} \\
\pv{managedBy} & \pv{PVSystem} & exactly one \pv{ControlAgent} \\
\pv{hasObservation} & \pv{PVSystem} & \pv{SystemObservation} \\
\pv{hasSignal} & \pv{PhotovoltaicArray} & \pv{TimeSeriesSignal} \\
\pv{hasRepresentation} & \pv{TimeSeriesSignal} & \pv{PatchedRepresentation} \\
\pv{assessesRepresentation} & \pv{ObservabilityAssessment} & exactly one \pv{PatchedRepresentation} \\
\pv{hasGenerationContext} & \pv{SystemObservation} & zero or one \pv{GenerationContext} value \\
\pv{hasAnomaly} & \pv{SystemObservation} & zero or more \pv{AnomalyState} values \\
\pv{hasObservabilityState} & \pv{ObservabilityAssessment} & one \pv{ObservabilityState} value \\
\pv{selectsMode} & \pv{PowerRoutingDecision} & exactly one \pv{ControlMode} value \\
\pv{makesDecision} & \pv{ControlAgent} & \pv{PowerRoutingDecision} \\
\pv{basedOn} & \pv{PowerRoutingDecision} & observation, assessment, or inferred state \\
\pv{hasRoute} & \pv{PowerRoutingDecision} & zero or one \pv{PowerRoute} value \\
\pv{hasSourceAsset} & \pv{PowerRoutingDecision} & zero or one \pv{EnergyAsset} \\
\pv{hasDestinationAsset} & \pv{PowerRoutingDecision} & zero or one \pv{EnergyAsset} \\
\bottomrule
\end{tabular}
\caption{Object-property signatures.}
\label{tab:objects}
\end{table}

\pv{hasRepresentation} is intentionally non-functional because one signal can be evaluated under several $(P,S)$ configurations. Each assessment identifies exactly which representation it evaluates. Each decision is based on exactly one current observation and its completed assessment; its own \pv{selectsMode} value therefore scopes the mode without accumulating incompatible modes on the persistent agent. The range of \pv{basedOn} is the union of observation and inferred-assessment entities, rather than the intersection that multiple RDFS range declarations would imply.

OWL's open-world semantics do not by themselves validate that endpoint types match a named route. The operational validator therefore admits only the seven ordered pairs in \autoref{eq:routes}, rejects identical endpoints, and checks consistency among \pv{hasRoute}, source, and destination. A hold decision has no route or endpoints and has a zero or absent setpoint. The \pv{basedOn} relation is retained only as the minimum evidence link required to justify a selected mode; no separate trace model is introduced.

\subsection{Data properties and units}

The synchronized observation combines the measured photovoltaic and environmental variables with asserted or simulated battery, load, and grid-reference values. The control-critical quantities are PV power $P_{\mathrm{pv}}$, tilted irradiance $G_{\mathrm{tilt}}$, normalized battery state of charge $s\in[0,1]$, and aggregate load demand $P_{\mathrm{load}}\geq0$. Each has a separate validity flag derived from presence, finiteness, freshness, and a configured physical range. Horizontal irradiance, air temperature, wind speed, and wind direction remain environmental evidence but do not independently block the controller. The signed \pv{gridExchangeReferenceW} is optional, with positive values denoting export; it is usable only when present, finite, and fresh.

Missing or stale data cannot be inferred safely from OWL's open-world assumption alone. An operational validation layer therefore compares the observation time stamp with a configurable staleness limit, validates numeric ranges, and asserts the four critical Boolean flags consumed by the ontology. This keeps data-quality validation distinct from semantic classification.

\begin{table}[H]
\centering
\small
\begin{tabular}{@{}L{0.29\linewidth}L{0.63\linewidth}@{}}
\toprule
\textbf{Owner} & \textbf{Data properties} \\
\midrule
\pv{SystemObservation} & \pv{pvPowerW}, \pv{powerRampRateWPerS}, \pv{loadDemandW}, optional \pv{gridExchangeReferenceW}, \pv{tiltedIrradianceWm2}, \pv{horizontalIrradianceWm2}, \pv{airTemperatureC}, \pv{windSpeedMs}, \pv{windDirectionDeg}, \pv{batterySoC}, \pv{observationTimestamp}, \pv{pvPowerValid}, \pv{tiltedIrradianceValid}, \pv{batterySoCValid}, and \pv{loadDemandValid}. \\
\pv{BatteryStorage} & \pv{capacityWh}, \pv{reserveSoC}, \pv{targetSoC}, \pv{maximumSoC}, \pv{maximumChargePowerW}, and \pv{maximumDischargePowerW}. \\
\pv{TimeSeriesSignal} & \pv{signalFrequencyHz}, \pv{samplingFrequencyHz}, \pv{periodSamples}, \pv{amplitudeW}, \pv{timeSpanSeconds}, and \pv{phaseOffsetSamples}. \\
\pv{PatchedRepresentation} & \pv{patchSize}, \pv{stride}, and derived \pv{overlapRatio}. \\
\pv{ObservabilityAssessment} & \pv{analysisStatus}, \pv{assessmentMethod}, \pv{assessmentMetric}, \pv{strideLockingRatio}, \pv{patchLockingRatio}, \pv{lockingTolerance}, \pv{effectEstimate}, \pv{posteriorLossProbability}, \pv{credibleIntervalLower}, \pv{credibleIntervalUpper}, and \pv{assessmentTimestamp}. \\
\pv{PowerRoutingDecision} & optional non-negative \pv{powerSetpointW}, \pv{decisionConfidence}, \pv{batteryCommand} in \{Charge, Discharge, Hold\}, \pv{pvCommand} in \{Connect, Isolate\}, and \pv{gridCommand} in \{Import, Export, Balance\}. \\
\bottomrule
\end{tabular}
\caption{Quantitative property groups and units encoded by their names.}
\label{tab:data}
\end{table}

A valid patched representation satisfies $P\geq1$ and $1\leq S\leq P$. Overlap is descriptive metadata:
\begin{equation}
\label{eq:overlap}
\gamma=\frac{P-S}{P}.
\end{equation}
It is not used as a standalone rule for declaring aliasing.

\section{Inference Rules}

\subsection{Environmental generation context}

Let $g_0<g_1<g_2$ be site-configurable tilted-irradiance thresholds and let $G_t$ denote the validated value of $G_{\mathrm{tilt}}$ at time $t$. The active generation context is
\begin{equation}
\label{eq:generation-context}
C_G(G_t)=
\begin{cases}
\valueof{Nighttime}, & G_t\leq g_0,\\
\valueof{Overcast}, & g_0<G_t<g_1,\\
\valueof{PartialCloud}, & g_1\leq G_t<g_2,\\
\valueof{ClearSky}, & G_t\geq g_2.
\end{cases}
\end{equation}
Values $g_0=0$, $g_1=150$, and $g_2=800\ \mathrm{W\,m^{-2}}$ are retained only as an illustrative configuration compatible with the current coursework. They are not universal meteorological boundaries. In particular, low irradiance can also be caused by solar position; the four values must therefore be read as generation-context proxies rather than certain weather diagnoses.

\subsection{Anomaly detection and precedence}

Let $v_P(t)$, $v_G(t)$, $v_s(t)$, and $v_L(t)$ denote the validity of PV power, tilted irradiance, battery SoC, and aggregate load demand. The control-critical validity predicate is
\begin{equation}
\label{eq:validity}
V_t=v_P(t)\wedge v_G(t)\wedge v_s(t)\wedge v_L(t).
\end{equation}
If $V_t$ is false, no physical anomaly or forecast-driven route is inferred from that observation. A valid irradiance value may still support a descriptive generation context, but the operational validator asserts \valueof{SensorMalfunction} and the cycle proceeds directly to safe-mode selection. Within a validated cycle, absence of the three listed anomaly values is interpreted locally as ``none detected''; this is a scoped closed-world policy, not a global OWL assumption.
The signed ramp rate, measured in watts per second when $\Delta t$ is expressed in seconds, is
\begin{equation}
\label{eq:signed-ramp}
r_t=\frac{P_t-P_{t-1}}{\Delta t}.
\end{equation}
The anomaly rules are expressed with configurable thresholds: $g_{\mathrm{trip}}$ for sufficient irradiance, $p_0$ for near-zero power, and $r_{\max}$ for maximum admissible ramp magnitude.
\begin{align}
\neg V_t
&\longrightarrow \valueof{SensorMalfunction},
\label{eq:sensor-malfunction}\\
V_t\wedge G_t\geq g_{\mathrm{trip}}
\wedge |P_t|\leq p_0
&\longrightarrow \valueof{InverterTrip},
\label{eq:inverter-trip}\\
V_t\wedge V_{t-1}\wedge |r_t|>r_{\max}
&\longrightarrow \valueof{RampEvent}.
\label{eq:ramp-event}
\end{align}
No numerical value is inherited for $r_{\max}$ because it must be calibrated after invalid values and outliers have been excluded. The controller evaluates sensor validity first. A sensor malfunction therefore cannot be reinterpreted as a physical ramp or inverter trip. If a valid inverter trip also satisfies the ramp predicate, the safe response to the trip takes precedence over ramp smoothing.

\subsection{Signal observability and structural aliasing}

For a signal frequency $f$, sampling frequency $f_s$, patch size $P$, and stride $S$, define a stride-lock ratio and a patch-periodicity ratio
\begin{equation}
\rho_S=\frac{fS}{f_s},
\qquad
\rho_P=\frac{fP}{f_s}.
\label{eq:locking-ratio}
\end{equation}
A stride phase-lock candidate $L_S$ and a within-patch repetition candidate $L_P$ are
\begin{equation}
\begin{aligned}
L_S&=\big[\exists m\in\mathbb{Z}^{+}:|\rho_S-m|\leq\varepsilon_S\big],\\
L_P&=\big[\exists k\in\mathbb{Z}^{+}:|\rho_P-k|\leq\varepsilon_P\big],\\
C&=(L_S\vee L_P)\wedge\left(0<f<\frac{f_s}{2}\right),
\end{aligned}
\label{eq:locking-candidate}
\end{equation}
where $\varepsilon_S$ and $\varepsilon_P$ are tied to frequency resolution and are never hardcoded frequency bands. Only $L_S$ expresses equality of consecutive patch phases; $L_P$ records whole-cycle repetition inside a patch and is retained as an H3 experimental factor. Their union $C$ is a candidate geometry class, not proof that two domain concepts are indistinguishable.

Let $q$ denote the task-relevant information being tested, for example frequency or amplitude. Let $\Delta_{\mathrm{loss}}(q)$ be its predeclared behavioral or representational loss relative to matched controls, $\delta$ the maximum practically negligible loss, $D$ the experimental evidence, and $\tau>1/2$ the posterior confidence threshold. Define confirmed information loss $B$ and confirmed preservation $R$ as
\begin{equation}
B=\left[\Pr\!\left(\Delta_{\mathrm{loss}}(q)>\delta\mid D\right)\geq\tau\right],
\qquad
R=\left[\Pr\!\left(\Delta_{\mathrm{loss}}(q)\leq\delta\mid D\right)\geq\tau\right].
\label{eq:bayesian-evidence}
\end{equation}
Because $\tau>1/2$, $B$ and $R$ cannot both hold for the same assessment. The assessment interface is designed to receive the planned offline evidence already specified in Deliverable~1: a local amplitude-recovery contrast with a heavy-tailed Student-$t_4$ likelihood and a Gamma regression with log link for non-negative prequential-MDL codelength. The population contrast uses the weakly informative prior $\bar\beta\sim\text{Student-}t_4(0,0.5)$; the Gamma model uses $\theta_{\mathrm{lock}}\sim\mathcal N(0,0.5)$ and shape $k\sim\mathrm{Gamma}(2,0.1)$. The stored record identifies the tested $(P,S)$ configuration, matched controls, method and metric, effect estimate, $95\%$ credible interval, and posterior probability of at least $20\%$ attenuation. These are evidence fields, not universal ontology thresholds. Until that planned inference has completed and satisfies the predeclared criterion, no instance may be asserted as \valueof{StructurallyAliased}.

The three observability values are assigned as follows:
\begin{equation}
C_O=
\begin{cases}
\valueof{StructurallyAliased}, & A,\\
\valueof{Observable}, & \neg A\wedge R,\\
\valueof{PartiallyObservable}, & \neg A\wedge\neg R,
\end{cases}
\qquad A=C\wedge B.
\label{eq:observability-state}
\end{equation}
\valueof{Observable} therefore means that positive evidence shows the tested representation preserves the task-relevant information within tolerance $\delta$. It does not mean global observability of the physical system, nor does it merely mean that aliasing was not found. Candidate geometry may still be \valueof{Observable} when the experiment demonstrates preservation. \valueof{PartiallyObservable} is the epistemic fallback: preservation has not been certified, while structural aliasing has not been confirmed. It includes a candidate with inconclusive evidence, an unfinished test, a credible interval crossing the decision threshold, conflicting metrics, or no preservation evidence. The term means partially trusted, not that a fixed fraction of samples or sensors is visible.

H1 compares each candidate frequency with matched nearby non-candidate controls and predicts higher MDL codelength or lower amplitude recovery; absent or non-localized loss refutes H1. H2 tests whether the loss at a candidate frequency remains approximately stable across sampled phase offsets; systematic phase dependence refutes that hypothesis. H3 tests whether candidate degradation sites move proportionally to $1/S$ or $1/P$ when the corresponding parameter varies. These hypotheses contribute controlled evidence to $D$, but remain falsifiable experimental factors rather than ontology subclasses or automatic action rules. Sensor validity is modeled separately as an anomaly and never relabels a model representation as structurally aliased.

When an experimental generator provides a phase $\phi$ in radians, the corresponding temporal offset is converted before storage:
\begin{equation}
o_{\mathrm{samples}}=\frac{\phi f_s}{2\pi f}.
\label{eq:phase-conversion}
\end{equation}
Neither $\phi$ nor $o_{\mathrm{samples}}$ is compared directly with stride to infer an anomaly.

The semantic propagation path is explicit:
\begin{equation}
\begin{aligned}
\pv{TimeSeriesSignal}
&\longrightarrow \pv{PatchedRepresentation}\\
&\longrightarrow \pv{ObservabilityAssessment}\\
&\longrightarrow \pv{ControlMode}
\longrightarrow \pv{PowerRoutingDecision}.
\end{aligned}
\label{eq:semantic-path}
\end{equation}
Aliasing can therefore reduce confidence in forecast-derived production, ramp expectations, forecast-driven anomaly recognition, and routing. It does not alter the directly observed values of $G_{\mathrm{tilt}}$, $P_{\mathrm{pv}}$, load demand, or SoC; those values are governed by the separate validity path.

\section{Closed-Loop Control}

\subsection{Mode selection}

The agent evaluates each synchronized control cycle using the strict priority
\begin{equation}
\label{eq:priority}
\valueof{SafeMode}\succ
\valueof{ConservativeMode}\succ
\mathrm{RampOverride}\succ
\mathrm{NormalRouting}.
\end{equation}
Let $A_{\mathrm{safe}}$ mean that \valueof{SensorMalfunction} or \valueof{InverterTrip} is present. The selected mode is
\begin{equation}
\label{eq:mode-selection}
M_t=
\begin{cases}
\valueof{SafeMode}, & A_{\mathrm{safe}}\vee
(C_O=\valueof{StructurallyAliased}),\\
\valueof{ConservativeMode}, & C_O=\valueof{PartiallyObservable},\\
\valueof{NormalMode}, & \text{otherwise}.
\end{cases}
\end{equation}
In \valueof{SafeMode}, the battery command is Hold and the grid balances the internal load. The PV command is Isolate only for an inverter trip; for sensor malfunction or confirmed structural aliasing it remains connected, but forecast- and ramp-driven battery decisions are suspended. When load demand itself is invalid, Balance carries no invented numerical setpoint. In \valueof{ConservativeMode}, the same seven route types remain available with a configurable reduced battery-power cap, normal target/reserve bounds, and low decision confidence. A concurrent ramp cannot activate either extended ramp bound; this implements the precedence in \autoref{eq:priority}.

\subsection{Routing policy}

The ordinary policy serves the internal load before charging the battery or exporting energy. Let
\begin{equation}
\label{eq:surplus-deficit}
P_{\mathrm{sur}}=(P_{\mathrm{pv}}-P_{\mathrm{load}})^{+},
\qquad
P_{\mathrm{def}}=(P_{\mathrm{load}}-P_{\mathrm{pv}})^{+}.
\end{equation}
Here daylight abbreviates \{\valueof{ClearSky}, \valueof{PartialCloud}, \valueof{Overcast}\}. A cycle issues a route only for a strictly positive transfer and may issue several co-mode decisions. All setpoints are clipped to SoC, capacity, and charge/discharge limits.

\autoref{fig:routing-policy} shows the priority scan and \autoref{tab:routing} states its boundary cases. The conservative branch rejoins the ordinary load-first router under its reduced cap. The ramp branch is available only in \valueof{NormalMode}.

\begin{figure}[H]
\centering
\begin{tikzpicture}[
	flow/.style={-{Latex[length=2mm]}, semithick},
	dashedflow/.style={-{Latex[length=2mm]}, semithick, dashed},
	start/.style={ellipse, draw, semithick, align=center, font=\scriptsize,
		text width=2.8cm, minimum height=8mm, inner sep=2pt},
	gate/.style={ellipse, draw, align=center, font=\scriptsize,
		text width=3.15cm, minimum height=10mm, inner sep=2pt},
	action/.style={draw, rounded corners=6pt, align=center, font=\scriptsize,
		text width=3.25cm, minimum height=11mm, inner sep=3pt},
	leaf/.style={draw, rounded corners=6pt, align=center, font=\scriptsize,
		text width=2.65cm, minimum height=19mm, inner sep=3pt},
	edgelabel/.style={font=\scriptsize, fill=white, inner sep=1pt}
]
\node[start] (cycle) at (0,0) {Synchronized control cycle};
\node[gate] (safeq) at (0,-1.45) {Safe trigger?\\Sensor fault, inverter trip, or structural aliasing};
\node[action] (safe) at (-5.0,-1.45) {\textbf{SafeMode}\\Battery Hold; Grid$\to$Load when observed. Trip: Isolate PV; otherwise PV remains connected.};
\node[gate] (partialq) at (0,-3.15) {$C_O=$ PartiallyObservable?};
\node[action] (conservative) at (5.0,-3.15) {\textbf{ConservativeMode}\\Low confidence; reduced battery cap; target/reserve bounds.};
\node[gate] (rampq) at (0,-4.85) {Valid RampEvent in NormalMode?};
\node[gate] (signq) at (-4.5,-6.55) {Ramp sign};
\node[gate] (ordinary) at (3.0,-6.55) {Ordinary load-first router\\Context, balance, and SoC};

\node[leaf] (positive) at (-6.4,-9.35) {\textbf{Positive ramp}\\PV$\to$Load; surplus PV$\to$Battery up to maximum; residual PV$\to$Grid.};
\node[leaf] (negative) at (-3.2,-9.35) {\textbf{Negative ramp}\\Load first; optional Battery$\to$Grid only toward a valid positive export reference.};
\node[leaf] (surplus) at (0,-9.35) {\textbf{Day surplus}\\PV$\to$Load; PV$\to$Battery to target; residual PV$\to$Grid.};
\node[leaf] (deficit) at (3.2,-9.35) {\textbf{Day deficit}\\PV$\to$Load; Battery$\to$Load to reserve; Grid$\to$Load residual.};
\node[leaf] (night) at (6.4,-9.35) {\textbf{Night}\\Battery/Grid$\to$Load; below reserve, Grid$\to$Load first, then Grid$\to$Battery.};

\draw[flow] (cycle) -- (safeq);
\draw[flow] (safeq.west) -- node[edgelabel, above] {yes} (safe.east);
\draw[flow] (safeq) -- node[edgelabel, right] {no} (partialq);
\draw[flow] (partialq.east) -- node[edgelabel, above] {yes} (conservative.west);
\draw[flow] (partialq) -- node[edgelabel, right] {no} (rampq);
\draw[flow] (rampq) -| node[edgelabel, near start, above] {yes} (signq);
\draw[flow] (rampq) -| node[edgelabel, near start, above] {no} (ordinary);
\draw[dashedflow] (conservative.south) |- node[edgelabel, near start, right] {ordinary bounds} (ordinary.east);
\draw[flow] (signq) -- node[edgelabel, left] {$r_t>r_{\max}$} (positive);
\draw[flow] (signq) -- node[edgelabel, right] {$r_t<-r_{\max}$} (negative);
\draw[flow] (ordinary) -- node[edgelabel, above, sloped] {surplus} (surplus);
\draw[flow] (ordinary) -- node[edgelabel, right] {deficit} (deficit);
\draw[flow] (ordinary) -- node[edgelabel, above, sloped] {night} (night);
\end{tikzpicture}
\caption{Load-first routing decision tree. Elliptical bubbles are ordered gates; rounded bubbles are terminal policies. Safe and conservative decisions precede the normal-mode ramp override, while the ordinary router covers daylight surplus, daylight deficit, and nighttime operation.}
\label{fig:routing-policy}
\end{figure}

\begin{table}[H]
\centering
\footnotesize
\begin{tabular}{@{}L{0.23\linewidth}L{0.29\linewidth}L{0.39\linewidth}@{}}
\toprule
\textbf{Condition} & \textbf{Ordered route set} & \textbf{Constraint} \\
\midrule
Daylight surplus & PV$\to$Load; PV$\to$Battery; PV$\to$Grid & Serve demand, charge to target, then export the residual. \\
Daylight deficit & PV$\to$Load; Battery$\to$Load; Grid$\to$Load & Discharge only to reserve; the grid covers the remaining deficit. \\
Night, $s>s_{\mathrm{reserve}}$ & Battery$\to$Load; Grid$\to$Load & Battery serves demand to reserve; the grid supplies the residual. \\
Night, $s\leq s_{\mathrm{reserve}}$ & Grid$\to$Load; if $s<s_{\mathrm{reserve}}$, Grid$\to$Battery & Supply the load first; charge only below reserve and stop at reserve. \\
Positive ramp with surplus & PV$\to$Load; PV$\to$Battery; PV$\to$Grid & NormalMode may charge beyond target, never beyond maximum SoC. \\
Negative ramp with valid positive reference & load-serving routes; Battery$\to$Grid & Export only after serving the load, within residual discharge capacity and above reserve. \\
Negative ramp without usable reference & ordinary surplus/deficit routes & Missing, stale, non-finite, zero, or negative reference disables export catch-up. \\
ConservativeMode & ordinary load-first routes & Reduced battery cap; target/reserve bounds; no ramp extensions. \\
SafeMode & Grid$\to$Load when observed & Battery Hold; grid Balance; isolate PV only for inverter trip. \\
\bottomrule
\end{tabular}
\caption{Routing matrix for the seven admissible transfer directions.}
\label{tab:routing}
\end{table}

For a negative ramp, Battery$\to$Grid is admissible only after the internal load has been fully served, the export reference is present, fresh, finite, and positive, measured export is below that reference, the battery remains above reserve, and discharge capacity remains after load support. If any condition fails, the ordinary load-first route applies. At night, Grid$\to$Battery begins only after Grid$\to$Load has been satisfied and stops when reserve is restored. These rules preserve \autoref{eq:power-balance} without turning the routing policy into a separate optimisation problem.

\printbibliography

\end{document}
